\textbf{(i)} $\Rightarrow$ If $fg=1$ then $a_0b_0=1$ and $a_0$ is a unit.

$\Leftarrow$ Suppose $a_0$ is a unit, WLOG say $a_0=1$. Then use the identity $(1-t)(1+t+t^2+\cdots)=1$ via substituting $t=-(a_1x+a_2x^2+\cdots)$.

\textbf{(ii)} Say $f^n=0$, then $a_0^n=0$. Consider the degree $n$ term, its coefficient is a sum of multiples of $a_0$ and $a_1^n$, and therefore $a_1^n$ is in the nilradical, and since nilradical is radical, $a_1$ is nilpotent. An induction with similar procedure shows the rest of the $a_i$ are all nilpotent.

Converse is not true. Take $A=k[x_1,x_2,\ldots]/(x_1^2,x_2^2,\ldots)$ and $a\in R$ non-zero nilpotent. Let $f=\sum x_i x^i$. Then $f^n$ always has term $x_1\cdots x_n x^{1+2+\cdots+n}$ not canceled out.

\textbf{(iii)} If $f$ in Jacobson radical, then $1+fg$ unit for all $g$. Take $g$ constants in $A$ and by (i), $a_0$ is in Jacobson radical. Conversely, if $1+a_0b$ unit for all $b$, then $1+fg$ unit for all $g$ by (i).

\textbf{(iv)} By (i), adding multiples of $x$ to a non-unit is non-unit. Thus $\mathfrak m+(x)$ is not $(1)$, and therefore by maximality $\mathfrak m+(x)=\mathfrak m$ and $x\in\mathfrak m$. Now the contraction, again by (i), is not $(1)$, so it is contained in some maximal ideal of $A$, adding this maximal ideal to $\mathfrak m$ we see we get an ideal that is not $(1)$, so it must be the contraction of $\mathfrak m$, and we can conclude $\mathfrak m$ is generated by $\mathfrak m^c$ and $x$.

\textbf{(v)} Consider the ideal generated by the prime ideal and $x$ which we see is prime.
